\chapter{The Framework of Nearly Everything}
\label{chap:framework-of-everything}

\section{The Ontology}

We are now in possession of all the pieces of our cosmological puzzle.

Our theory begins with minimal assumptions.

\subsection*{The Transcendence of Information}

If reality is framed in terms of information, a natural question arises: Is the total information of the cosmos finite or infinite?

If it were finite, one would have to ask: Finite by how many bits? And who, or what, set that arbitrary threshold?
Any fixed binary bound would itself require an external explanation, pushing the mystery of creation back one level.
The only non-arbitrary conclusion is that the deep nature of everything cannot be finitely bounded; it must be infinite, existing beyond any baseline description.

Pushing this logic to its necessary conclusion shows that not even information itself can be fundamental. 

To posit a binary bit as an ontological primitive is to assume a universe carved into discrete, bounded choices (e.g., $0$ or $1$).
Yet, if a discrete digital state is possible, a perfectly smooth, continuous, infinite-state wave is equally possible.
Declaring the baseline of reality to be inherently digital is an arbitrary axiomatic choice that itself demands a prior mechanism to enforce that discreteness. 

Therefore, we do not presuppose that reality is fundamentally discrete, informational, or governed by prior, external laws.
The raw totality ($\mathbf{U}$) is completely unformatted and transcendent. 

Information is not an inherent property of objective reality, but an emergent artifact of observation.
Based on all observational evidence, we as human beings possess a strictly finite informational structure.
We cannot ingest an infinite, continuous continuum in its raw form.
To exist as stable, bounded entities, we must parse the infinite through a discrete, interpretative filter.
The ``bit'' is not the fundamental building block of the cosmos; it is the fundamental unit of distinction employed by a finite interpreter.

\subsection*{The Guitar String Analogy}

An intuitive analog is a guitar string.
The string itself is perfectly continuous. It isn't made of discrete "tone particles."
Yet, when you pluck it, it can only vibrate at discrete, quantized frequencies (the fundamental note, the octave, the third harmonic, etc.), 
because the string is clamped at both ends! The boundaries force a continuous medium to manifest discrete states.

Discreteness does not exist as an intrinsic property of the physical substrate.
Rather, discreteness emerges solely when a finite length string (observer) asks: \emph{Which of these perfectly smooth waves is simpler?}

We are forced to ask this question because, under a Solomonoff or algorithmic measure, we overwhelmingly find ourselves embedded within the simplest,
most compressible configurations. Waves that settle into discrete, resonant, low-complexity modes achieve minimal description length and maximum probability. 



\section{The Initial Assumption: Configuration Space}

The only firm starting point we adopt is purely observational: conscious finite observers exist and have structured experiences.
For formal descriptive purposes, we introduce a finite, static configuration space:
\[
\mathcal{C} = \{0,1\}^n
\]
where each $c \in \mathcal{C}$ is a complete, static snapshot of bits. 

Each configuration $c$ admits multiple potential descriptions $d \in \mathcal{D}(c)$.
We quantify the total complexity $C_{\text{total}}(d)$ of a description by splitting it into its dual spectral and geometric representations, matching our informational Lagrangian:
\[
C_{\text{total}}(d) = C_s(d) + C_g(d)
\]

The spectral complexity $C_s(d)$ measures the algorithmic bit-cost required to instantiate and track the wave-like modes:
\[
C_s(d) = \sum_i \left[ \text{Cost}(A_i) + \text{Cost}(\phi_i) + \left(\frac{\omega_i}{\Delta \omega}\right) \right]
\]
where the sum runs over active spectral modes, $\text{Cost}(A_i)$ represents the bit-depth of the amplitude envelope, $\text{Cost}(\phi_i)$ is the fixed-width register encoding the periodic phase interval $[0, 2\pi]$, and $\omega_i / \Delta \omega$ represents the linear resource allocation required to track the mode frequency relative to the minimum resolution threshold $\Delta\omega$.

The geometric complexity term $C_g(d)$ accounts for the algorithmic overhead of localized spatial structures, boundaries, and metric fields.
Consistent with our physical foundations, the spectral cost scales linearly with the number of active modes multiplied by their frequency, ensuring an exact algorithmic analogue to thermodynamic energy.



\section{The Equation of Joint Compressibility}

We have established that the quantum mechanical wavefunction compresses the possibility space of our universe through a spectral language ($C_s$), while General Relativity compresses its realization space through a geometric language ($C_g$). 

Because the relationship between these two frameworks is a many-to-many mapping, any given observer wavefunction can be interpreted through nearly infinitely many different geometric spacetime metrics, and conversely, any spacetime manifold can be decomposed into nearly infinitely many different wavefunctions.
This creates an unimaginably massive configuration space. 

How does nature choose our specific reality from this vast haystack?

In information theory, when a system must be simultaneously validated across two dual representations, its joint probability is the product of its individual probabilities.
Therefore, the absolute measure of any emergent physical state ($s$) within the totality is given by the \textbf{Joint Compressibility Equation}:

\begin{equation}
P_{\text{joint}}(s) = P_{\text{spectral}}(s) \times P_{\text{geometric}}(s) = 2^{-C_s(s)} \times 2^{-C_g(s)} = 2^{-\left(C_s(s) + C_g(s)\right)}
\end{equation}

This deceptively simple formula yields a profound variational principle for a unified framework.
Nature minimizes the total description length.
The cosmos is governed by a singular, overarching informational Lagrangian:

\begin{equation}
\label{eq:lagrangian}
\mathcal{L}_{\text{info}} = C_s(s) + C_g(s) \longrightarrow \text{Minimum}
\end{equation}

Neither geometry nor waves are primary; they are equal, co-dependent projections of a single underlying informational matrix.
Our observed reality is the exact mathematical sweet spot of this joint optimization.
Spacetime is smooth and the quantum world is lawful because this specific configuration represents the absolute shortest combined code nature can write.



\section{Induced Measure and Observer Functionals}

The raw algorithmic weight of a static configuration $c$ is accumulated across all its valid descriptions:
\[
W(c) = \sum_{d \in \mathcal{D}(c)} 2^{-C_{\text{total}}(d)}
\]
This defines our raw baseline probability distribution over the state space:
\[
P(c) = \frac{W(c)}{\sum_{c'\in\mathcal{C}} W(c')}
\]

Because $C_s(d)$ maps complexity linearly to frequency, the calculation $2^{-C_{\text{total}}(d)}$ inside the summation guarantees that high-frequency states suffer genuine exponential probability suppression, mirroring the Boltzmann distribution ($P \propto e^{-\beta E}$) found in statistical mechanics.

An observer is modeled as a grading semantic functional:
\[
\mu_O : \mathcal{C} \to [0,1]
\]
which measures how strongly a given bit configuration realizes a stable, observer-relative structure.
This functional induces an observer-relative equivalence relation $c_1 \sim_O c_2$ whenever $c_1$ and $c_2$ realize the same relational observer structure to within the tolerances encoded by $\mu_O$. 

The observer-conditioned probability then becomes:
\[
P(c \mid O) = \frac{\mu_O(c)W(c)}{\sum_{c'} \mu_O(c')W(c')}
\]



\section{The Mathematical Unification}

To transition from the probability of a single static configuration $P(c \mid O)$ to the probability of an entire experienced history,
we must track how an observer traverses the configuration space.
Let a path $\gamma \in \Gamma_O$ represent a sequence of configurations ordered by our internal temporal relation $\prec_O$.
Because the configurations along the path are independent slices of the static totality, the total probability of the path is the product of its individual constituent measures, meaning their algorithmic costs sum additively.

By changing our informational units from bits to natural units (where $2^{-C} = e^{-C \ln 2}$), we can express the global probability of an entire observed history $\gamma$:
\[
\mathbb{P}(\gamma \mid O) = \frac{1}{Z_O} \exp\left(-\lambda\mathcal{C}_O[\gamma]\right), \qquad \gamma \in \Gamma_O
\]
where $\mathcal{C}_O[\gamma] = \int_{\gamma} C_{\text{total}}(c)\, dt$ represents the accumulated informational cost along the path, and $\lambda$ acts as an algorithmic inverse temperature (scaling the bit-to-nat conversion and the observer's filtering acuity).
This states that observed physical laws are simply the large-deviation minimizers of the total informational cost function $\mathcal{C}_O[\gamma]$ over the set of all semantic paths.




\section{The Emergence of Spacetime and Laws}

Our core hypothesis is that configurations capable of supporting observers strongly favor smooth, low-frequency spectral content.
Sharp, chaotic particle trajectories require massive high-frequency modes, which exponentially spike the linear description length $C_s(d)$ and are instantly suppressed by the joint measure. 

As a result, observer-compatible realities are overwhelmingly dominated by configurations exhibiting inertia, smooth curved paths, and law-like behavior.
Spacetime geometry emerges precisely because it makes macroscopic informational boundaries computationally cheap to describe ($C_g$). 

Time in this static block-universe is not a fundamental background dimension.
It emerges entirely as an internal ordering relation $\prec_O$ on configurations, where $c_1 \prec_O c_2$ if $c_2$ supports a more advanced, stable realization of the observer functional $\mu_O$.
The experienced flow of time corresponds to traversal along the most compressible, optimized paths through this timeless superposition.



